\subsection{The coefficient rings and sizes on the selected cycles}
\label{sec:host-edge-tests}
\hypertarget{source-host-edge-tests}{}

The preceding transfer and averaging estimates are conditional on local size bounds. We now check those bounds on the good cycles supplied by the arithmetic choice of primes.

Fix the target character and \(0<\epsilon<1/52\). Fix a candidate prime \(p\) of even Frobenius order and a good cycle as in Section~\ref{subsec:frobenius-selection}. For its Frobenius step \(a_p\), the cycle length satisfies
\[
 n'=d_p=\operatorname{ord}(a_p),\qquad 2\mid n'.
\] Write \(N\) for the
number of tops, \(r_G=K+1\) for the fixed plane dimension,
\(G'=G\) for the length bound there, and \(Q=p^{G'}\).
Numbering from the chosen orbit representative identifies its
\(s\)-th place with the translate by \(s a_p\). Thus the top
positions, the signs by parity and the two endpoint weights agree
with those in the local test.

At a top the first and last entries
are \(a_s,-b_s\). The other entries there and all entries at the
remaining places have absolute value at most \(\lfloor c_G\rfloor\).
Here
\[
 c_G=4S_c=16k\Delta_{\rm host},
\]
where \(\Delta_{\rm host}\)
is the size parameter bounding the middle weights. These bounds specify the weights used in the local lattice estimate.

\paragraph{The local field.}
Write \(W=W(\mathbb F_{p^{n'}})\) and \(K_0=\operatorname{Frac}W\).
Take
\[
 O_s=W(\mathbb F_{p^{2n'}})[w]
\]
with \(w^2=p\). Fix a
place \(\mathfrak w\mid p\) of the number field of coefficients
\(F\). Embed \(F\) in an algebraic closure of
\(K_s=\operatorname{Frac}O_s\) by embedding its completion
\(F_{\mathfrak w}\) there over \(\mathbb Q_p\). The field
\[
 K'=K_s F
\]
in that closure is finite over \(K_s\), since a
primitive element of \(F/\mathbb Q\) is algebraic over \(K_s\).

The ring \(O'\) of integers of \(K'\) is the integral closure
of \(O_s\), a finite free module over the complete DVR \(O_s\).
Its valuation, normalized by \(v(p)=1\), restricts to the chosen
valuation on \(F\). This is the ring used for the place in the
host construction. The maps on \(W,K_0\) are the composites of the
standard embeddings with \(O_s\to O'\), \(K_s\to K'\), and the
square root of \(p\) is the image of the specified \(w\).

\begin{lemma}[Finite base change of an edge test]\label{lem:edge-test-host-base-change}
For \(O_s\subset O'\) at the chosen place of the coefficient extension, keep the
same vertices over \(W\) and the same competing edges.
Write \(o,c\) for the order and conductor of a depth-\(D\) test with \(m\) successful directions over \(O_s\). Extension to \(O'\) preserves these four parameters
\[
 (m',D',o',c')=(m,D,o,c).
\]
 The ambient dimension is unchanged.
The valuation normalization is \(v(p)=1\) in both rings.
\end{lemma}

\begin{proof}
\proofstep{The source and target over the actual local rings.}
The vertices are still lattices over \(W\) in \(K_0^k\).
In particular the same edges and all their competitors are
present after enlarging coefficients. At position \(s\) the
twisted lattice is spanned by the images of a \(W\)-basis under
the \(s\)-th Frobenius followed by the embedding. This span,
and the sum of the two spans with multipliers \(1,w\) at an
edge, commute with extension to \(O'\). One may evaluate the
Weyl lattice in any one frame of this sum. A second frame
differs by an integral invertible change of basis and acts by
an integral automorphism, including its determinant twist.

The products of flag minors generating the lattice in this
frame extend coefficient by coefficient. Thus for a lattice \(H\) in one frame its realization after extension of coefficients is
\[
 H_{O'}=\sum_{h\in H}O'\,\iota(h)\simeq O'\otimes_{O_s}H,
\]
where \(\iota\) denotes coefficient extension in that frame.
The same assertion holds for their tensor product on a cycle.
The tableau basis gives the same dimension over both fraction
fields. In the dual space extend each linear functional in that
basis. Integrality on the old lattice then implies integrality
on the new span. Multiplying the endpoint and edge lattices by
their defining powers of \(p,w\) commutes with the maps.

For a target evaluation map with chosen source vectors and
target covectors giving \(p^D I_m\), take their images at the
new place. If \(\iota:O_s\to O'\) is the coefficient map, the evaluation on the chosen test vectors has matrix
\[
 \bigl(\iota(p^D\delta_{ij})\bigr)_{i,j}=p^D I_m.
\]
This follows by extending the bilinear evaluation on each frame basis. The new
lattice conditions at the same vertices are given by those
basis spans, so extending each vector and covector preserves
them. In particular the preimages of the scaled coordinate
vectors still give a right inverse after the same scaling. All required memberships at the
other edges are preserved. The number \(m\) and its denominator
\(\dim V\) are identical. This proves finite base change even
when \(O'/O_s\) is ramified.  \label{cor:host-edge-tests-1}

\end{proof}

\paragraph{Counting tops and bounding the parameters of the auxiliary construction.}
Summing condition (iii) defining a good cycle over its \(n'\)
possible interval starts counts each top \(G'\) times. Also every
gap between consecutive even tops is at most \(G'\). Thus
\[
 n'/G'\le N\le n'N_1/G'\le 40n'\delta=:N_+ .
\]
For the last step use
\(N_1\le2+\max(6eG'\delta,(h+3)\log_2 x)\), with
\(G'\delta\ge (6r_G+1)(h+3)10^{6r_G+1}\log x\).
In particular \(N\gg_\epsilon x^{\epsilon/8}/\log x\),
as \(n'\ge x^{\epsilon/4}\), and \(N_+\le20N_*\) with
\(N_*=2n\delta\). The same marked blocks give at least two
segments and their length and interior counts are bounded by
\(G',N_1\).

We have \(3G'\le n'\) and
\(2L_1(p)\le L_2<\lceil1/\delta\rceil\).
In both ranges use
\[
 \Lambda=C_0 N_1^2\log a,\quad a=x\exp(x^\epsilon),
 \quad b'=a/X',\quad X'=\exp(x^{\epsilon/3}).
\]
Here and below the constants in the local estimates are fixed
before \(x\). From the rank bounds of
Section~\ref{subsec:frobenius-selection}, \(n\ll x^{2+8\epsilon}\),
\(x^{9/2+10\epsilon-\epsilon/100}\le k\le2x^{9/2+10\epsilon}\)
and \(N_1=O_\epsilon(\log x)\). The discriminant bound
\(\log\Delta_{\rm host}\le x^{\epsilon/4}/2\) gives
\(\log(1+c_G)\ll x^{\epsilon/4}+\log x\).
All comparisons below are uniform in the candidate prime, the
cycle and the weights within the indicated intervals.

\paragraph{Large primes.}
Put \(Y=x^{13},\ Y_b=x^\epsilon/20\) and let \(I_b\)
denote the length \(N_* x^\epsilon\Lambda b'/k\) of the
second weight interval. Then
\[
 k^2Y^2\le p<X',\qquad
 Y_bN\Lambda b'/k\le I_b\le b'/2,\qquad
 2Y_bN\Lambda+C N\Lambda^2=o(k)
\]
for each fixed \(C\). Indeed \(N\le40n\delta\) and both
\(x^\epsilon\Lambda\) and \(\Lambda^2\) are
\(x^{2\epsilon+o(1)}\).
We have \[
 \begin{gathered}
 2N_1+2\le k-1, \\ 8\log(NY)\le\Lambda, \\ c_G\le\Lambda b'/k, \\ kY^2Y_bN\Lambda\log a\le X',
\end{gathered}
\] the last left side being
polynomial in \(x\). These verify the spacing and width
conditions even at the ends of the integer intervals.

\paragraph{The hypotheses for the large-prime test.}
To check that the weight ranges remain separated, use the width parameter
\(\Delta=C_{\Delta}N_1\mathcal L/k+2(b_h-b_l)/b'\),
where the local reference \(a\) may be rounded up by one and
\(\mathcal L\) is then computed there. The original real interval
has width \(I_b\). A shift by \(\lfloor c_G\rfloor\), the lower
boundary allowance and the enlarged box add \(O(\Lambda b'/k)\)
to the widths on each side of this interval
(the shift from the harmonic decomposition is bounded by
\(\Lambda b'/k\)). Thus a band of relative half-width
\[
          4I_b/b'+C_\rho N_1\Lambda/k
\]
covers the local band of half-width \(2\Delta\), taking the
fixed \(C_\rho\) large enough for these shifts and rounding
the center upward by at most \(1/b'\).

The possible increase
relative to the center \(a/b'\) is \(O(1/a)\).
This is the band included in the exceptional set of
Section~\ref{subsec:frobenius-selection}. We may take
\(C_\Delta\) large enough that the enlarged weight deviation
is at most \(\Delta/2-18N_1\mathcal L/k\) and
\(40N_1\mathcal L/k\le\Delta\). We have \(\Delta<1/6\)
eventually by the bound on the exceptional width.

\paragraph{Depth and conductor on the selected Frobenius cycle.}
The depth for this cycle is
\[
 D_{\rm I}=\left\lceil C_D N_+Y^2(N_1^5+k^3)\log_p a\right\rceil .
\]
\begin{proposition}[Comparison of the selected lattices for large primes]
\label{prop:host-large-prime-test}
In the stated large-prime range, the rank estimate applies
on the selected lattices at depth \(D_{\rm I}\) and conductor
\(8(D_{\rm I}+1)\). If \(\bar\delta_{\rm I}\) denotes the resulting mean deficit, the depth, conductor and deficit satisfy
\[
 D=D_{\rm I},\qquad c=8(D_{\rm I}+1),\qquad\bar\delta_{\rm I}=O(NN_1/x^{13}).
\]

\end{proposition}

\begin{proof}
The defining inequality makes \(D_{\rm I}\) dominate the local
depth, and the size bounds make it polynomial in \(x\).
Since \(G'\ll_\epsilon x^{\epsilon/8}\log x\) and \(p<X'\),
\(\log Q\ll_\epsilon x^{11\epsilon/24}\log x=o(x^\epsilon)\).
Consequently
\[
 C k^3NQ^4\big(8(D_{\rm I}+1)+(N+n')c_G\big)Y=o(b')
\]
for every fixed \(C\), by taking logarithms. The larger internal
conductor in Section~\ref{sec:edge-transfer}, which includes
\(O(Nk\log_p a)\), satisfies the same inequality.

In particular
\(kQ\le N_1\mathcal L b'/k\). Dividing the expression for \(E_0\)
in Section~\ref{sec:symbol-correction} by \(a\), using
\(p\ge k^2Y^2\), \(b'/a=1/X'\) and the same bound on
\(Q^2c/a\), gives \((E_0+Q)/a=o(N_1\mathcal L/k)\).
The large-prime estimate therefore applies at the prescribed depth
and conductor \(8(D_{\rm I}+1)\). From \(p<a/b'\) we obtain both comparisons needed in the order estimate,
\[
 \frac{b'}p<\frac a{p^2},\qquad \frac{\mathcal L b'}k<\frac{\mathcal L a}{pk}.
\]

The local deficit bound is unchanged by the finite scalar
extension in Lemma~\ref{lem:edge-test-host-base-change}: both
the test rank and its denominator are dimensions of the same
extended vector spaces. Thus it gives \(O(NN_1/x^{13})\)
on the host.
\end{proof}

\paragraph{Smaller primes.}
Proposition~\ref{prop:residue-plane-neighbor-test} gives the
tests in this range. We check its size conditions on the
same cycles and level choices as in the larger-prime case.

\begin{proposition}[Comparison of the selected lattices for smaller primes]
\label{prop:host-small-prime-test}
For \(13\le p<p_0=k^2x^{26}\) on the same selected cycles,
the residue-plane comparison applies at order
\(\lceil t_{\rm II}\rceil\) and conductor \(8(D_{\rm II}+1)\),
with the depth and order in \eqref{eq:small-prime-depth-order}.
Writing \(\kappa\) for its deficit, the resolution and the bound are
\[
 o=\lceil t_{\rm II}\rceil,\qquad c=8(D_{\rm II}+1),\qquad\kappa=O(Nx^{-12}).
\]

\end{proposition}

\begin{proof}
Set \(Y_2=x^{26}\),
\(c_0=2b'+2c_G\), \(A=a-c_G\), and
\(L_1=\lceil\log_p((kx^{26})^4)\rceil\). The interval conditions
give \(A+\lfloor c_G\rfloor\le a_s\le 2A-\lfloor c_G\rfloor\)
and \(b_s+2\lfloor c_G\rfloor\le c_0\), as \(c_G=o(b')\).
For large \(x\) the lower bound on \(N\) gives
\(N_1^2\le r_G N\). Also \(2\le r_G\le Y_2\),
\(N_1+3r_G\le k\), and \(C N_1^2\log a\le k\). The spacing
condition follows from
\[
 p^{L_1-1}\ge k^4 x^{104}/p
 \ge k^2 x^{78}\ge256k^2 N Y_2^2 .
\]
Here the last inequality uses \(N\le n\le x^3\).
The prescribed depth and real order are
\begin{equation}\label{eq:small-prime-depth-order}\begin{aligned}
 D_{\rm II}={}&\left\lceil C_D' N_+(N_1^3+r_G^5+k^2)\log_p a
                 +2r_Gn' c_G+r_G N_+c_0\right\rceil,\\
 t_{\rm II}={}&C_t r_G N\big(\Lambda a/(pk)+a p^{-L_1}
                         +c_0+kQ\big)+2r_G n'c_G .
\end{aligned}\end{equation}

Now \(\log Q\ll_\epsilon x^{\epsilon/8}(\log x)^2
=o(x^{\epsilon/3})\), since \(p\) is polynomially bounded.
Thus \(D_{\rm II}\le x^{O(1)} b'\) and
\[
 C k^3r_G^2NQ^5\big(8(D_{\rm II}+1)+r_G N c_0+
                 2r_G n'c_G\big)Y_2=o(a)
\]
by \(a/b'=X'\). These are all the numerical hypotheses for
the smaller-prime test at the host. Section~\ref{sec:plane-lifts}
proves the estimate on edges with residue quotient of dimension
\(r_G\), at order \(\lceil t_{\rm II}\rceil\) and conductor
\(8(D_{\rm II}+1)\). With the bounds on \(r_G,k\) and
\(Y_2=x^{26}\), the deficit
\(\min(1,CNr_G^2 k/Y_2)\) is \(O(Nx^{-12})\).
Lemma~\ref{lem:edge-test-host-base-change} preserves both
ranks in this fraction. \label{cor:host-edge-tests-2}
\end{proof}
